% Transcription — fonds Alexandre Grothendieck, shelfmark 86, batch 2 (pages 21-32).
% Established from the facsimile archives/batches/86/batch-02.pdf, cut from a
% copy of 86.pdf fetched through the project's own relay
% (https://grothendieck-relay.onrender.com/source/86.pdf), not uploaded by the
% user; PDF page k = archivists' page 20+k, no cover sheet. Per the skill
% transcribe-grothendieck. The folder is thirty-two pages; this is the second
% and last batch. Batch 1 (pages 1-20) was being transcribed in parallel by
% another pass in parallel; it was committed before this pass closed and was
% read then. It breaks off on p. 20 at « induisent », which p. 21 continues
% (« des permutations paires sur l'ens des sommets ... »). Its notation
% agrees with this batch (\underline{a}, \mathfrak{P}_2, Pent, Rep, tr); its
% G_s^+ is written G_s here, as on the leaves. The one disagreement, its
% « trièdre » against this pass's first reading « tridron », was settled on
% the leaf: re-cropped at 900 dpi on pp. 23, 25 and 26, the word is « tri »,
% a dot standing for the accent, then « dres » — the same letter shapes as
% « icosaèdre » two lines above it on p. 26. It is read « trièdre »
% throughout, with one \note at the first occurrence. The labels
% Corollaire, Proposition, Th, Dém, NB are set in bold, as in batch 1. This
% batch has no heading of its own until his « Reconstruction ... » on p. 26:
% batch 1's section « Cours --- Polyèdres réguliers III » (a cover's words)
% runs on into it.
%
% Dating and title. The dating is the folder's own, copied verbatim from the
% catalogue; the group « Géométrie et topologie combinatoire » (69-89) holds
% it, and since the folder has a date of its own the group's range is not
% added. The inventory title « [Polyèdres réguliers] » is in square brackets:
% it was supplied by the archivists, not written by Grothendieck, and is kept
% in \foldertitle with its brackets and never used as a heading of his. No
% leaf of this batch carries a title of his own.
%
% Pages skipped: none. Pages summarised: none (nothing typed, nothing by
% another hand). Pages transcribed from his hand: 21-32. Pages 21-30 are one
% continuous run in bright blue felt-tip ink, fast but mostly legible; the
% lower half of page 28 is a palimpsest (a cancelled attempt at b), struck
% line by line and overwritten) and is given as one \struck{\ill{}} with a
% note. Pages 31-32 are in a different, dark blue-black ink and a steadier
% hand: a table and three formulas, the rest of page 32 blank.
%
% His pagination: none visible on any leaf. Page 21 opens in mid-sentence
% (« des permutations paires sur l'ens des sommets ... »), continuing a
% page of batch 1; the object Π, the sets S, A, F, T and the notation tr for
% the « trièdre distingué » of an edge direction were therefore introduced
% there.
%
% What the batch is about. The combinatorial icosahedron Π and the classical
% isomorphism of its rotation group with the alternating group on five
% letters, obtained intrinsically through the five-element set T of its
% « trièdres distingués ». P. 21: G_s (stabiliser of a vertex) acts on T by
% even permutations, because every vertex s defines, via P(s) ≃ T, a
% pentagonal structure π_s on T, hence an orientation ω_s; for a combinatorial
% polygon with n ≡ 1 (4) sides the set of vertices is canonically oriented.
% A repère r = {s, a, f} gives generators σ0, σ1, σ2 of G = Aut(Π). P. 22
% (with a drawing of the icosahedron, the fundamental triangle shaded and
% σ0, σ1, σ2 marked): σ1, σ2 fix s, σ0 fixes the third vertex w of f, so all
% act evenly; Corollary G → Alt_T; since ω_{u(s)} = u(ω_s) and G is
% transitive on S, Corollary: the ω_s are all equal, a canonical orientation
% of T preserved by G. P. 23: Proposition on a face f = {u, v, w} and the
% three adjacent faces: f induces the circular permutation on two triples of
% edge directions, corresponding to the two orientations of f; the three
% distinguished trihedra tr of the edges of f and t_ω, t_ω' are distinct;
% Corollary T ≃ f ⊔ Or(f). P. 24: compatibility with G_f ≃ D_3 = 𝔖_3, so
% G_f acts faithfully on T with image a dihedral subgroup of order 6, hence
% the image of G in Alt_T has order divisible by 6, by 10, so by 30;
% Proposition for an edge a = {s, t}: the five elements tr{s,u} (s ∈ a,
% u ∈ φ) and tr ā of T are distinct. P. 25: T ≃ {tr ā} ⊔ a × φ, G_a ≃
% 𝔖_a × 𝔖_φ ≃ Z/2 × Z/2 acts faithfully, so the image has order divisible
% by 4, hence 60 = |Alt_T|; Theorem: G = Aut Π → 𝔖_{T(Π)} has image Alt_T
% and kernel the subgroup 𝔷 of order 2 generated by the antipodism; G ≃
% 𝔷 × G⁺. P. 26: Corollary G⁺ ≃ Alt_T; « Reconstruction de l'icosaèdre
% orienté par l'ens orienté T »: of the ½·4! = 12 pentagonal structures on T,
% 6 belong to each orientation; S/a ≃ Pent⁺(T, ϖ). P. 27: the set 𝔸(T) of
% triples (t, γ, α) (γ a square structure on T \ {t}, α one of its two
% pairs of opposite edges), of cardinal 5·½·3!·2 = 30, splits into two
% packets of 15 by orientation; A/a ≃ 𝔸⁺(T, ϖ), F/a ≃ 𝔓₂(T). P. 28: the
% dictionary s̄ ↔ π, ā ↔ (t, γ, α), f̄ ↔ Ω for the « icosaèdre gauche »
% Π̄ = Π/a, and a Proposition on the incidence relations: a) s̄, ā incident
% iff γ is the square deduced from π at t, with α = {{u,v},{ũ,ṽ}}; b)
% (mostly cancelled) ā, f̄ incident iff t ∈ T \ Ω. P. 29: N.B. on the two
% square structures and the choice of γ by orientation; R' = A^↑(Π̄) ≃ pairs
% of disjoint Ω, Ω'; c) s̄, f̄ incident iff Ω is an anti-edge (diagonal) of
% π; d) repères of Π̄ ↔ triples (π, τ, Ω). P. 30: a repère defines a
% numbering T = {t0, ..., t4} compatible with ϖ and conversely; left to do:
% make S̃, Ã, F̃ explicit in terms of (T, ϖ), and Π̄ as a quotient polyhedron.
% Pp. 31-32 (the other ink): a diagram of the flag-type sets D^0, D^1, D^2,
% D^{01}, D^{12}, D^{02}, D^{012} of the icosaèdre gauche and the table of
% their descriptions in terms of a five-element set E with an orientation ω
% (S ≃ Pent_ω(E), A ≃ Car(E), F ≃ Tr_3(E), A⃗ ≃ Pent^s_ω(E), A^↑,
% A^∧ ≃ Pent^a_ω(E) ≃ Pent^s_ω(E), A^{↑→} ≃ Rep_ω(E) total orders
% compatible with ω, S̃, Ã, F̃ as circular permutations, oriented squares,
% oriented triangles); the generators σ0, σ1, σ2 as explicit permutations of
% (a1, ..., a5), and ρ_f = σ1σ0, ρ_s = σ2σ1, σ = σ0σ1.
%
% His underlinings in the prose are rendered \emph (the subset has no text
% underline); \underline is kept only inside mathematics.
%
% Notation fixed once for the batch: Π the icosahedron, S, A, F its sets of
% vertices, edges, faces; T its set of distinguished trihedra and tr the map
% from edge directions to T (the overline on {u,v} marks the direction, as on
% the leaves); G = Aut(Π), G_s, G_f, G_a stabilisers, G⁺ the orientation-
% preserving subgroup; Alt_T and 𝔖_T for the alternating and symmetric groups
% of T (he writes « Alt » in Roman and a gothic S), 𝔖_3, 𝔖_a, 𝔖_φ likewise;
% D_3 dihedral; the canonical orientation of T as \varpi (he writes an omega
% with a bar over it, as in 75); ω_s the orientations defined by the
% vertices; 𝔓₂(T), 𝔓₃(S) as \mathfrak{P}; his double-struck A as
% \mathbb{A}; the antipodism as \underline{a} (an underlined a, also in the
% quotients S/a, A/a, F/a); the curly figure-3 naming the kernel of order 2
% as \mathfrak{z}, per hand.md; Pent⁺, Pol, Ar, Car, Tr, Rep, Or kept as he
% abbreviates them; the icosaèdre gauche Π̄ and its sets S̄, Ā, F̄ with bars.
% On pp. 31-32 the flag superscripts of A are rendered A^{\uparrow}, A^{\wedge}
% and A^{\uparrow\to} (the last is an up-arrow joined at its foot to a
% right-arrow on the leaf).
%
% Readings to check first: the opening of p. 21 (the struck words, « ou,
% ... », « type n ≡ 1 (4) », « canoniquement », « On verra plus bas »); the
% end of the σ0 sentence on p. 22 (« ... une \ill{} paire de T »); on p. 23
% the three overlined edge directions of a) (letters and subscripts), the
% inserted « dist. » and the struck word before « : », and the first of the
% three tr-equalities of b); on p. 25 the subscript of G (a or s), the
% phrase between « Comme » and « et \underline{a} renverse », and the target
% of G⁺ ≃ (read 𝔄_5, doubtful); on p. 27 the superscript of 𝔸'(T), the
% margin « = Ens des 2 diagonales du ... », and the openings « \ill{}, on a »
% and « Enfin, \ill{} »; the whole of b) on p. 28; on p. 29 the parenthesis
% after « adjacents » in d); on p. 30 the symbol before (T, π); on p. 31 the
% overwritten symbols Car, Tr_3 and the word read « pointés ».
%
% Pass: Opus 5.5 (claude-opus-5-5), 2026-09-26 — first pass, unchecked against the pages by a human.

\documentclass[11pt,a4paper]{article}
% Le préambule commun à toutes les transcriptions du fonds Grothendieck.
%
% Il tient en une idée : **l'incertitude se compose, elle ne se résout pas.**
% Une transcription qui lisse un mot douteux a détruit la seule chose pour
% laquelle on la faisait — un lecteur doit pouvoir savoir, à chaque endroit,
% ce qui était sur le feuillet et ce qui a été ajouté. Les six macros qui
% suivent sont donc l'appareil critique, et non de la décoration.
%
% Les mêmes commandes sont reconnues par `scripts/render.mjs`, qui produit la
% vue de lecture du site. Une macro ajoutée ici doit l'être aussi là-bas,
% faute de quoi le rendu échouera bruyamment — ce qui est voulu : un rendu
% silencieusement faux est pire qu'un rendu absent.

\ProvidesPackage{grothendieck}[2026/08/08 Transcriptions du fonds Grothendieck]

\usepackage{amsmath,amssymb,amsthm}
\usepackage{mathrsfs}
\usepackage{stmaryrd}
% Les diagrammes commutatifs sont partout dans ce fonds : c'est la langue
% même de la géométrie algébrique de Grothendieck, et une transcription qui
% les rendrait en prose ne transcrirait rien.
\usepackage{tikz-cd}
% Français par défaut — c'est la langue des feuillets — mais l'anglais est
% chargé aussi : la traduction et le résumé sont des documents anglais, et la
% typographie française (espaces avant les deux-points) y serait une faute.
% Ces documents basculent par \selectlanguage{english} en tête de corps.
\usepackage[english,french]{babel}
\usepackage{fontspec}
\usepackage{geometry}
\geometry{a4paper,margin=25mm,marginparwidth=28mm}
\usepackage{marginnote}
\usepackage{xcolor}
% ulem, not soul. `soul`'s \st and \ul are built on a character-by-character
% scan that XeLaTeX's font machinery breaks: the document fails with an
% undefined control sequence rather than degrading. ulem does the same two jobs
% and is Unicode-engine safe. `normalem` stops it hijacking \emph, which the
% transcriptions use for Grothendieck's own emphasis.
\usepackage[normalem]{ulem}
\usepackage{hyperref}
\hypersetup{colorlinks=true,linkcolor=black,urlcolor=black,pdfborder={0 0 0}}

% --- Métadonnées du lot -----------------------------------------------------
% Reprises telles quelles par le script de rendu, qui les affiche en tête de
% la vue de lecture. Elles fixent ce que le fichier prétend couvrir : sans
% elles, une transcription flotte sans point d'attache dans un fonds de seize
% mille feuillets.

\newcommand{\folder}[1]{\gdef\grothFolder{#1}}
\newcommand{\batch}[1]{\gdef\grothBatch{#1}}
\newcommand{\leaves}[2]{\gdef\grothFirst{#1}\gdef\grothLast{#2}}
\newcommand{\foldertitle}[1]{\gdef\grothFoldertitle{#1}}
\gdef\grothFolder{?}\gdef\grothBatch{?}\gdef\grothFirst{?}\gdef\grothLast{?}\gdef\grothFoldertitle{}

% --- Le résumé --------------------------------------------------------------
%
% Il ouvre la lecture modernisée, sous le titre « Résumé », et il est composé
% en petit. Ce n'est pas un ornement : c'est le seul passage du document qui
% ne soit pas des mathématiques, et un lecteur doit pouvoir le reconnaître
% avant de l'avoir lu — puis le sauter s'il connaît déjà le sujet, ou s'y
% arrêter s'il ne le connaît pas. Le filet qui le suit marque la reprise du
% corps.

\newenvironment{resume}
  {\small\setlength{\parskip}{0.45em}}
  {\par\medskip\hrule\medskip}

% --- Mots-clés ---------------------------------------------------------------
%
% En anglais, en clôture du résumé de la lecture modernisée : le vocabulaire
% moderne sous lequel chercher ce que les feuillets construisent. Le manifeste
% du site (`npm run manifest`) les extrait pour étiqueter la cote entière —
% cette ligne est la source unique des tags, il n'y a pas de fichier à côté.
\newcommand{\keywords}[1]{\par\smallskip\noindent
  {\footnotesize\textsc{Keywords} --- \textit{#1}}\par}

% --- L'appareil critique ----------------------------------------------------

% Le numéro de feuillet, en marge. C'est l'ancre qui permet de régler un
% désaccord sur une formule en regardant *un* feuillet plutôt qu'une chemise
% de six cents. Le site s'en sert aussi pour tourner les pages du fac-similé
% quand on fait défiler la transcription.
\newcommand{\leaf}[1]{%
  \marginnote{\footnotesize\color{gray!70}\textsf{#1}}%
  \label{leaf:#1}\ignorespaces}

% Illisible : on ne devine pas. Un mot inventé qui se lit comme les autres est
% le pire résultat possible.
\newcommand{\ill}{\textcolor{red!60!black}{[\,\ldots\,]}}

% Lecture douteuse : le mot est donné, et signalé comme incertain.
\newcommand{\uncertain}[1]{\uline{#1}}

% Ajout de l'éditeur — une lettre manquante, un mot sous-entendu.
\newcommand{\add}[1]{\textcolor{blue!50!black}{[#1]}}

% Biffé par Grothendieck lui-même. Ce qu'il a rayé fait partie du document :
% une rature dit souvent où la pensée a changé de direction.
\newcommand{\struck}[1]{\sout{#1}}

% Note du transcripteur, sur l'état matériel du feuillet.
\newcommand{\note}[1]{\par\smallskip{\small\color{gray!60!black}\textit{[#1]}}\par\smallskip}

% Note marginale de Grothendieck, distincte du corps du texte.
\newcommand{\marginal}[1]{\par\smallskip\noindent\hspace{1em}%
  {\small\color{gray!60!black}#1}\par\smallskip}

% --- Titre ------------------------------------------------------------------

\AtBeginDocument{%
  \begin{center}
    {\large\bfseries Fonds Alexandre Grothendieck --- cote n\textsuperscript{o}~\grothFolder}\\[2pt]
    {\itshape\grothFoldertitle}\\[4pt]
    {\small feuillets \grothFirst--\grothLast\ (lot \grothBatch)}\\[2pt]
    {\footnotesize\color{gray!60!black}Transcription établie d'après le fac-similé numérisé
     par l'Université de Montpellier.\\
     Les lectures incertaines sont soulignées, les illisibles notées [\,\ldots\,],
     les ajouts entre crochets.}
  \end{center}
  \vspace{1em}}

\endinput


\folder{86}
\batch{2}
\pages{21}{32}
\dating{[à partir de 1977]}
\watermark{Édition de démonstration}
\foldertitle{[Polyèdres réguliers] : notes manuscrites (s.d.)}

\begin{document}

\page{21}
\note{la page commence au milieu d'une phrase, suite d'un feuillet du lot 1}

\struck{par \ill{}} des permutations paires sur l'ens des sommets (ou,
\ill{}, sur l'ens des côtés). \struck{Donc} On en conclut que $G_s$ opère
\add{fidèlement}\note{« fidèlement » est inséré au-dessus de la ligne} sur
$T$ par permutations paires. Aussi, on trouve que pour tout polygone
combinatoire de type \uncertain{$n \equiv 1 \ (4)$}, l'ens des sommets est
\uncertain{canoniquement} orienté, via épinglage qui permet de numéroter
l'ens des sommets (\uncertain{et} l'orientation de l'ens des sommets ne
dépend pas de l'épinglage…).

On peut aussi dire que tout $s \in S$ définit sur $T$, via
$P(s) \xrightarrow{\sim} T$, une structure polygonale (pentagonale)
$\pi_s$\note{« $\pi_s$ » est écrit au-dessus de la ligne, rattaché par un
trait à « structure »}, qui est invariante sous $G_s$ (qui opère \ill{}
\ill{}) \ill{} groupe des automorphismes de $\pi_s$. A fortiori, $G_s$
invarie l'orientation correspondante de $T$, soit $\omega_s$, donc $G_s$
opère par permutations de $T$ paires. [On \uncertain{verra plus bas} que
$\omega_s$ ne dépend pas en fait du choix de $s \in S$]

Fixons-nous un repère $r = \{s, a, f\}$ de $\Pi$, d'où des automorphismes
$\sigma_0, \sigma_1, \sigma_2$ engendrant $G = \operatorname{Aut}(\Pi)$.

\page{22}
\note{à gauche, un dessin de l'icosaèdre en perspective, arêtes cachées en
pointillé ; sommets marqués $u_1$, $w_1'$, $u'$, $v_1'$, $w$, $v$, $v_1$,
$w_1$, $u = s$, $u_1'$, $w'$, $v'$ ; un triangle fondamental hachuré, de
sommet $s$, porte sur ses côtés les marques $\sigma_0$, $\sigma_1$,
$\sigma_2$}

$\sigma_1$ et $\sigma_2$ fixent $s$, \add{donc} \struck{conservent
l'orientation de} \add{induisent des permutations}\note{la correction est
écrite au-dessus de la ligne biffée} paires de $T$. D'autre part
$\sigma_0$ conserve le troisième sommet $w$ (opposé à $a$) de la face
(triangulaire) $f$, donc par ce qui précède \struck{conserve}
\add{\uncertain{induit}} \ill{} une \ill{} paire de $T$. Donc

\textbf{Corollaire} $G = \operatorname{Aut}(\Pi)$ opère sur $T$ par
permutations paires, i.e.\ \ill{}
\[
  G \longrightarrow \operatorname{Alt}_T
\]

Comme $u \in G$, $s \in S$ implique évidemment
\[
  \omega_{u(s)} = u(\omega_s)
\]
et puisque $G$ est transitif sur $S$

\textbf{Corollaire} Les orientations $\omega_s$ ($s \in S$) de $T$
associées aux $s \in S$ sont toutes égales\note{« associées aux $s \in
S$ » est rejeté au début de la ligne suivante, précédé d'un trait
vertical}. D'où une \emph{orientation} \emph{canonique sur $T$}
(évidemment conservée par $G = \operatorname{Aut}(\Pi)$).

\page{23}

\textbf{Proposition} Considérons une face $f = \{u, v, w\}$.
\struck{Alors les trois arêtes} Soit $t_f = t$ l'ens des trois arêtes de
$f$ (en correspondance 1-1 avec l'ens de ses sommets), et soient $\{u_1,
v_1, w_1\}$ les troisièmes sommets des \add{trois} faces \add{$\neq
f$}\note{« trois » et « $\neq f$ » sont insérés au-dessus de la ligne}
adjacentes aux arêtes. Alors

a) $f$ induit la permutation circulaire \struck{\ill{}} sur les
\add{directions} d'arêtes $\overline{\{u, v\}}$, $\overline{\{v,
w_1\}}$, $\overline{\{w, u_1\}}$ et sur l'ens des directions d'arêtes
$\overline{\{u, w_1\}}$, $\overline{\{w, v_1\}}$, $\overline{\{v, u_1\}}$,
correspondant \add{\uncertain{dist.}} aux trièdres\note{le mot s'écrit « tri » suivi d'un point pour l'accent, puis « dres », exactement comme « icosaèdre » à la page 26 ; on lit « trièdres », comme le lot 1} \struck{\uncertain{opposés}}
: ces orbites de $\{\omega, \omega^{-1}\} = \{\omega, \omega'\}$ \ill{},
en corr.\ 1-1 avec l'une des deux orientations (i.e.\ des 2 ordres
circulaires) sur $f$, sont $\{t_\omega, t_{\omega'}\}$.\note{lecture douteuse des lettres et des indices de ces trois
directions}

\marginal{$uvw$ \quad $uwv$}\note{dans la marge gauche, en regard des deux
triples de directions : les deux ordres circulaires sur $f$}

b) Les trois trièdres distingués $\operatorname{tr}\overline{\{v, w\}} =
t_u$, $\operatorname{tr}
\overline{\{w, u\}} = t_v$, $\operatorname{tr}\overline{\{u, v\}} = t_w$,
$t_\omega$ et $t_{\omega'}$ sont distincts.\note{la première égalité est lue avec doute}

\textbf{Dém} a) est clair. \struck{Pour} $t_u, t_v, t_w$ \uncertain{deux
à deux} distincts est ce qu'on peut \ill{} de la prop.\ précédente, ainsi
que $t_u \neq t_\omega = \operatorname{tr}\overline{\{w, u_1\}}$, $t_u \neq
t_{\omega'} = \operatorname{tr}\overline{\{v, u_1\}}$ et de même pour
$t_v$, $t_w$.

\textbf{Corollaire} On a une bijection canonique
\[
  T \simeq f \amalg \operatorname{Or}(f)
\]

\page{24}

(où $f$ est identifié à un élément de $\mathfrak{P}_3(S)$).

Cette bijection est évidemment compatible avec l'opération de $G_f \simeq
D_3$ $= \mathfrak{S}_3$, opérant sur \struck{\ill{}} $\operatorname{Or}(f)$ via le
quotient $\mathfrak{S}_2$ de $\mathfrak{S}_3$ (i.e.\ via signature) et sur
$f$ tautologiquement. \struck{Donc} [On retrouve que $G_f$
\struck{conserve} \add{opère} par permutations paires sur $T$ ---] On voit
ainsi que \struck{l'image de} $G_f$ \struck{dans} \add{opère fidèlement
dans} $T$, et son image dans $\mathfrak{S}_T$ (contenue dans
$\operatorname{Alt}_T$) est donc un s-gpe diédral $6$.
\emph{Donc} l'image de $G$ dans $\operatorname{Alt}_T$ est d'ordre un
multiple de $6$, et par les résultats précédents un multiple de $10$, donc
un multiple de $30$.\note{« $= \mathfrak{S}_3$ » est inséré au-dessus de
la ligne}

\textbf{Proposition} Soit $a = \{s, t\} \in A$. Considérons
\struck{l'ens $\varphi$ des deux} \add{les deux} faces incidentes à $a$, et
\struck{l'ens $\varphi$ $\{s, \ldots\}$} \add{l'ens $\varphi$ de} leurs
troisièmes sommets $u$, $v$, \struck{dans $A$ les arêtes} Considérons
\[
  a \times \varphi \longrightarrow A
\]
défini par $(s, u) \longmapsto \{s, u\}$. Alors les cinq éléments
$\operatorname{tr}\overline{\{s, u\}}$ ($s \in a$, $u \in \varphi$),
$\operatorname{tr}\bar{a}$ de $T$ sont distincts.

Résulte des deux prop.\ précédentes.

\page{25}

On trouve ainsi une bijection canonique
\[
  T \simeq \{\operatorname{tr}\bar{a}\} \amalg a \times \varphi ,
\]
évidemment compatible avec l'action de $G_a$ $\simeq \mathfrak{S}_a \times
\mathfrak{S}_\varphi \simeq \mathbf{Z}/2\mathbf{Z} \times
\mathbf{Z}/2\mathbf{Z}$, qui opère donc \emph{fidèlement} sur $T$. On
retrouve que $G_a$ opère sur $T$ par permutations paires, et de plus on
conclut que l'image de $G$ dans $\operatorname{Alt}_T$ est d'ordre
divisible par $4$. Et comme déjà \add{d'ordre} divisible par $30$, il est
divisible par $60$. Or $\operatorname{Alt}_T$ est d'ordre $60$,
\struck{donc} et \struck{Aut} $G$ d'ordre $120$, on trouve\note{l'indice est douteux :
$a$ ou $s$ ; de même trois lignes plus bas}

\textbf{Th} L'homom.\ canonique
\[
  \underset{\text{icosaèdre}}{G = \operatorname{Aut}\Pi}
  \longrightarrow
  \mathfrak{S}_{\underset{\text{ens des trièdres distingués}}{T(\Pi)}}
\]
a comme image $\operatorname{Alt}_T$, et comme noyau le sous-groupe
$\mathfrak{z}$\note{la lettre en forme de 3 bouclé est insérée au-dessus de
la ligne ; on la rend par $\mathfrak{z}$, comme ailleurs dans le fonds}
d'ordre $2$ engendré par l'antipodisme $\underline{a}$.

Comme \struck{\ill{}} $\Pi$ \ill{} et $\underline{a}$ renverse
l'orientation, on trouve\note{un « Corollaire » biffé en tête de ces lignes}
\[
  G \simeq \mathfrak{z} \times G^{+}
\]
donc
\[
  G^{+} \xrightarrow{\sim} \mathfrak{A}_5
\]
\note{la lettre après la flèche est lue $\mathfrak{A}$ avec doute}
donc

\page{26}

\textbf{Corollaire} $G^{+} \xrightarrow{\sim} \operatorname{Alt}_T$

\section*{Reconstruction de l'icosaèdre \emph{orienté} par l'ens orienté $T$}
\note{titre de sa main, marqué d'un trait vertical dans la marge, non
souligné en entier}

Soit $\Pi$ icosaèdre, \struck{\ill{}} pas orienté pour le moment. L'ens $T$
de ses trièdres \emph{distingués} est orienté de façon canonique, soit
$\varpi$ son orientation. Sur $T$ il y a exactement $\frac12 4!$ $(=12)$
structures polygonales, qui se partagent en $2$ paquets de $6$ structures
polygonales, associées à l'une et l'autre orientation. Soit
$\operatorname{Pent}^{+}(T, \varpi)$ l'ens des $6$ structures pentagonales
associées à $\varpi$. On a une application canonique
\[
  S/\underline{a} \longrightarrow \operatorname{Pent}^{+}(T, \varpi)
\]
compatible avec op.\ de $G$. Or $G$ opère transitivement sur $S$, et il est
immédiat que $\operatorname{Alt}_T$ opère transitivement sur
$\operatorname{Pent}^{+}(T, \varpi)$ (vrai pour tout polygone à $n \equiv 1
\ (4)$ côtés…), donc l'application précédente est surjective, donc
bijective
\[
  S/\underline{a} \xrightarrow{\sim} \operatorname{Pent}^{+}(T, \varpi) .
\]

\page{27}

Soit d'autre part
\[
  \mathbb{A}'(T) = \left\{ (t, \gamma, \alpha) \;\middle|\;
  \begin{array}{l}
    t \in T \\
    \gamma \in \operatorname{Pol}(T \smallsetminus \{t\}) \\
    \alpha \in \operatorname{Ar}(\gamma)/\text{antipod.}
  \end{array}
  \right\}
\]
\note{un premier symbole biffé devant $\mathbb{A}$ ; l'exposant de
$\mathbb{A}$, lu comme un accent, disparaît dans la suite, où il écrit
$\mathbb{A}(T)$. Dans le triple, la deuxième lettre est surchargée ($\pi$
corrigé en $\gamma$ ?) ; les conditions sont écrites
$\pi \in \operatorname{Pol}$, $\alpha \in \operatorname{Ar}(\pi)$ sans
correction}
\marginal{$= $ Ens des $2$ diagonales du \ill{}}\note{sous
« $/\text{antipod.}$ », dans la marge droite}

On voit que le cardinal de $\mathbb{A}'(T)$ est
\[
  5 \cdot \tfrac12 (4 - 1)! \cdot 2 = 30
\]
\struck{\ill{}}. Mais si $\alpha \in \mathbb{A}(T)$, $\alpha$ détermine une
bijection
\[
  T \simeq \{t\} \amalg a \times \varphi
  \qquad (\text{où } \varphi = \operatorname{Ar}(\pi) \smallsetminus a),
\]
d'où une orientation de $T$. Ainsi $\mathbb{A}(T)$ se décompose en $2$
paquets de $15$ éléments, associés aux deux orientations ($\varpi$,
$-\varpi$) de $T$ ; on voit tout de suite que sur chacun, telle
$\mathbb{A}^{+}(T, \varpi)$, le groupe alterné $\operatorname{Alt}_T$
opère transitivement.

\ill{}, on a une application canonique
\[
  A/\underline{a} \longrightarrow \mathbb{A}^{+}(T, \varpi) ,
\]
compatible avec $G = \operatorname{Aut}(\Pi)$, et on voit comme
précédemment qu'elle est surjective, donc bijective.

Enfin, \ill{} une application canonique
\[
  F/\underline{a} \longrightarrow \mathfrak{P}_2(T)
\]
compatible aux op.\ de $G$, donc (comme $\operatorname{Alt}_T$ est
transitif sur $\mathfrak{P}_2(T)$) surjective donc bijective.

\page{28}

Donc on a reconstruit les ens de sommets, d'arêtes, de faces de
l'\emph{icosaèdre gauche} en termes de $(T, \varpi)$.

Il faudrait reconstruire aussi les relations d'incidence.
\begin{align*}
  \bar{s} \in \bar{S} \quad &\longleftrightarrow \quad \pi \in
    \operatorname{Pent}^{+}(T, \varpi) \\
  \bar{a} \in \bar{A} \quad &\longleftrightarrow \quad \{t, \gamma, \alpha\}
    \in \mathbb{A}^{+}(T, \varpi) \\
  \bar{f} \in \bar{F} \quad &\longleftrightarrow \quad \Omega \in
    \mathfrak{P}_2(T)
\end{align*}
\note{dans la deuxième ligne, la lettre du milieu est surchargée ; dans la
troisième, une première écriture entre accolades est noircie}

\textbf{Proposition}

a) $\bar{s}$ et $\bar{a}$ incidents ssi $\gamma$ est la structure carrée sur
$T \smallsetminus \{t\}$ se déduisant de la structure pentagonale $\pi$ sur
$T$ en prenant les sommets adjacents $u$, $v$ de $t$, $\tilde{u} \neq t$
($\tilde{v} \neq t$) adjacents à $u$, [$v$], et la structure carrée de
diagonales $\{u, \tilde{u}\}$, $\{v, \tilde{v}\}$, donc \struck{\ill{}}
\[
  \begin{array}{c}
    u \text{---} v \\
    \tilde{v} \text{---} \tilde{u}
  \end{array},
  \qquad \text{et} \quad \alpha = \bigl\{ \{u, v\}, \{\tilde{u}, \tilde{v}\}
  \bigr\}
\]
[$\bar{R} = \vec{A}(\bar{\Pi})$ s'identifie à l'ens des couples $(\pi,
t)$, avec $\pi \in \operatorname{Pent}^{+}(T, \varpi)$, $t \in T$, donc
\[
  \bar{R} \simeq \vec{A}(\bar{\Pi}) \simeq S(\bar{\Pi}) \times T(\bar{\Pi})
  \quad ]
\]

b) $\bar{a}$ et $\bar{f}$ incidents ssi $t \in T \smallsetminus \Omega$,
\struck{\ill{}}

\note{le reste de b) est un palimpseste : quatre ou cinq lignes biffées
une à une puis surchargées, où l'on distingue une bijection
$\overline{T \smallsetminus \{t\}} \simeq \{\tau\} \amalg \ldots \amalg
\Omega$, « la condition que », « ($\gamma$) … tels que », « $\Omega$
sont des arêtes consécutives » et « donc $\alpha$ est déterminé par » ; rien
n'y est lisible de façon suivie}

\note{dans la marge gauche, trois croquis : un pentagone de sommets $t = \tau$,
$u$, $v$, $\tilde{u}$, $\tilde{v}$, de centre $s$, marqué $(\pi)$ ; sous
lui, une colonne de paires fléchées en cycle, $\{s, u\} \to \{s, v\} \to
\{s, \tilde{u}\} \sim \{t, v\}$, $\{s, \tilde{v}\} \sim \{t, u\}$, au-dessus
d'un premier essai biffé ; plus bas, un fragment d'icosaèdre de sommets $u$,
$v$, $w$, $u_1$, $v_1$, $w_1$, l'arête $\{u, v\}$ épaissie, et le cycle
$\{u, w_1\} \to \{w_1, v\} \to \{v, w\} \to \{w, u\}$}

\page{29}

\textbf{NB} Il y a précisément deux structures carrées sur $T
\smallsetminus \{\tau\}$\note{le $T$ est surchargé d'une tache} satisfaisant
à \ill{}, définies par les deux bijections possibles entre $\Omega
\smallsetminus \{\tau\} = \Omega'$ et $\Omega$ (relation de
non-adjacence), mais on constate que les $2$ \struck{signatures}
\add{orientations} de $T$ associées aux deux choix de $\gamma$, lorsque
$\alpha$ satisfait \ill{}, sont opposées, donc il y a exactement une telle
$\gamma$ telle que
\[
  (\tau, \gamma, \alpha) \in \mathbb{A}^{+}(T, \varpi) \ \ldots
\]
[$R' = A^{\uparrow}(\bar{\Pi}) \simeq \bigl\{ (\Omega, \Omega') \in
\mathfrak{P}_2(T) \times \mathfrak{P}_2(T) \bigm| \Omega \cap \Omega' =
\emptyset \bigr\}$]

c) $\bar{s}$ et $\bar{f}$ incidents ssi $\Omega$ \emph{anti-incidente}
pour la structure pentagonale $\pi$. [$R'' \simeq$ ens des structures
pentagonales sur $T$ munies d'une (anti) arête \add{au choix}]\note{« au
choix » est écrit sous « (anti) arête », avec une flèche}

\marginal{\textbf{NB} $a \sim \{u, v\}$}\note{sous un croquis de
pentagone à diagonales, de sommets $u$, $v$, marqué $a$, $s$, $f$}

d) $\{\bar{s}, \bar{a}, \bar{f}\}$ forment un repère ssi $\pi$, $(\tau,
\gamma, \alpha)$ et $\Omega$ se déduisent de
\[
  \{\pi, \tau, \Omega\}, \qquad \pi \in \operatorname{Pent}^{+}(T, \varpi),
  \quad \tau \in T, \quad \Omega \in \mathfrak{P}_2(T)
\]
satisfaisant\note{sous les trois lettres, trois petits signes
d'appartenance renvoient à $\operatorname{Pent}^{+}(T, \varpi)$, $T$,
$\mathfrak{P}_2(T)$} : $\tau \notin \Omega$, $\Omega$ anti-arête pour
$\pi$, $\tau$ \uncertain{adjacent} (\ill{} un seul sommet \ill{}) à $\Omega$
[pour $\pi$] [ce qui donne, pour \struck{\ill{}} $\pi$ ($6$ choix) et
$\tau$ ($5$ choix) donnés, $2$ choix pour $\Omega$], en prenant pour
$\gamma$ l'unique structure de carré sur $T \smallsetminus \{\tau\}$ pour
laquelle $\Omega$, $\{u, v\}$ (les deux sommets adj.\ à $t$) et $\{\bar{u},
\bar{v}\} = T \smallsetminus \{\tau\} \smallsetminus \{u, v\}$ sont des
arêtes, et $\alpha = \bigl\{ \{u, v\}, \{\tilde{u}, \tilde{v}\} \bigr\}$
\ill{} est la \ill{} orientation de $(T \smallsetminus \{\tau\}, \gamma)$
\ill{} par $\Omega$. La donnée

\note{dans la marge gauche, de haut en bas : un pentagone à diagonales,
sommets $u$, $v$, marqué $a$, $s$, $f$ ; une petite esquisse du diagramme
de la page 31 ($A^{\uparrow\to}$ en haut, $\vec{A}$, $\tilde{A}$,
$A^{\uparrow}$ au milieu, $S$, $A$, $F$ en bas), dont les flèches ne se
laissent pas toutes suivre ; un hexagone
entièrement biffé et, à côté, un croquis biffé ; un pentagone de sommets
$t$, $u$, $v$, $\tilde{u}$, $\tilde{v}$, de centre $s$ ; un carré de
sommets $u$, $v$, $\tilde{u}$, $\tilde{v}$, le côté $\{u, \tilde{v}\}$
épaissi et marqué $\Omega$}

\page{30}

des repères (définis par $\pi$, $\tau$, $\Omega$) définit un ordre sur
$T$, car $\pi$ définit une structure pentagonale, $t_0 = \tau$ une
origine pour celle-ci, et $\Omega$, diagonale partant de $t_1 =$ l'unique
élément $\in \Omega$ adjacent à $\tau$, on trouve un repère de
\ill{}$(T, \pi)$, d'où une numérotation
\[
  T = \{t_0, t_1, t_2, t_3, t_4\} ,
\]
\add{qui est compatible avec l'orientation $\varpi$}.\note{écrit
au-dessus de la ligne, sur une première insertion biffée, illisible}

La donnée d'une telle numérotation définit une structure pentagonale $\pi
\in \operatorname{Pent}^{+}(T, \varpi)$, \struck{$\Omega = \{t_1$} un
$\tau = t_0 \in T$, un
\[
  \Omega = \{t_1, t_3\} \in \mathfrak{P}_2(T) ,
\]
définissant un élément de $\operatorname{Rep}(\bar{\Pi})$ …

\note{dans la marge gauche, un pentagone, le sommet du bas marqué $\tau =
t_0$, celui de gauche $u = t_1$, la diagonale $\Omega$ en pointillé}

Il reste : expliciter [pour $\Pi$ orienté]
\[
  \left\{
  \begin{array}{l}
    \tilde{S}(\bar{\Pi}) \simeq S(\Pi) \\
    \tilde{A}(\bar{\Pi}) \simeq A(\Pi) \\
    \tilde{F}(\bar{\Pi}) \simeq F(\Pi)
  \end{array}
  \right.
\]
en termes de $(T, \varpi)$. Et expliciter aussi, \struck{\ill{}} de façon
\uncertain{directe}, la structure de $\bar{\Pi}$ comme polyèdre
quotient…

\page{31}
\note{les pages 31 et 32 sont d'une autre encre, bleu-noir, et d'une main
plus posée}


\begin{tikzcd}[column sep=small, row sep=small, nodes={font=\scriptsize}]
    & & A^{\uparrow\to} = D^{012} \arrow[dll] \arrow[dl] \arrow[d] \arrow[dr] \arrow[drr] & & \\
    \tilde{S} \arrow[d] & \vec{A} = D^{01} \arrow[dl] \arrow[dr] & \tilde{A} \arrow[d] & A^{\uparrow} = D^{12} \arrow[dl] \arrow[dr] & \tilde{F} \arrow[d] \\
    S = D^{0} & & A = D^{1} & & F = D^{2}
\end{tikzcd}



\begin{tikzcd}[column sep=small, row sep=small, nodes={font=\scriptsize}]
    & A^{\uparrow\to} = D^{012} \arrow[d] & \\
    & A^{\wedge} = D^{02} \arrow[dl] \arrow[dr] & \\
    S = D^{0} & & F = D^{2}
\end{tikzcd}

\note{les deux diagrammes sont côte à côte en tête de page ; les $D$ sont
écrits sous les lettres, reliés par un signe $=$ vertical}

\emph{Icosaèdre gauche} (sommets $S$ etc.…) en termes de $(E,
\omega)$, $E$ ens.\ à $5$ éléments, $\omega$ une orientation
\begin{align*}
  S &\simeq \operatorname{Pent}_\omega(E) && \text{struct.\ pent.\ comp.\ avec } \omega \\
  A &\simeq \operatorname{Car}(E) && \text{partitions du type } (2, 2, 1) \\
    & && \simeq \text{structures carrées sur une partie à 4 él.\ de } E \\
  F &\simeq \operatorname{Tr}_3(E) && \text{triangles} \subset E \\
  \vec{A} &\simeq \operatorname{Pent}^{s}_\omega(E) && \omega\text{-structures pentagonales à somm.\ marqué} \\
  A^{\uparrow} &\simeq \operatorname{Tr}_3(E) && \text{triangles pointés} \\
  A^{\wedge} &\simeq \operatorname{Pent}^{a}_\omega(E) \simeq \operatorname{Pent}^{s}_\omega(E)
    && \text{structures pol.\ à arête marquée} \\
  A^{\uparrow\to} &\simeq \operatorname{Rep}_\omega(E) && \text{ordres totaux sur } E \text{ comp.\ avec } \omega \\
  \tilde{S} &\simeq \overrightarrow{\operatorname{Pent}}_\omega(E) && \text{permutations circulaires comp.\ avec } \omega \\
  \tilde{A} &\simeq \overrightarrow{\operatorname{Car}}(E) && \text{carrés orientés} \subset E \\
  \tilde{F} &\simeq \overrightarrow{\operatorname{Tr}}_3(E) && \text{triangles orientés} \subset E
\end{align*}
\note{les symboles de droite sont souvent écrits par-dessus un premier
symbole : $\operatorname{Car}$ sur une première lettre illisible,
$\operatorname{Tr}_3$ sur un $\mathfrak{P}_3$ (?) pour $F$, $A^{\uparrow}$
et $\tilde{F}$, $\overrightarrow{\operatorname{Pent}}$ sur un premier mot
illisible. Les commentaires biffés : pour $F$ « \struck{partitions à 3
él.} », pour $\vec{A}$ « pol » corrigé en « pentagonales », pour
$A^{\uparrow}$ « \struck{parties à trois él.} » sous « triangles », pour
$\tilde{F}$ « \struck{parties à trois él.\ avec ordre circ.\ dessus} ».
Pour $A^{\uparrow}$, « pointés » est lu avec doute. Pour $A$, la ligne « partitions du type $(2,2,1)$ » est coupée à droite par
les croquis}

\[
  \left\{
  \begin{aligned}
    \sigma_0(a_1, a_2, a_3, a_4, a_5) &= (a_1, a_4, a_5, a_2, a_3) \\
    \sigma_1(a_1, a_2, a_3, a_4, a_5) &= (a_2, a_1, a_5, a_4, a_3) \\
    \sigma_2(a_1, a_2, a_3, a_4, a_5) &= (a_1, a_5, a_4, a_3, a_2)
  \end{aligned}
  \right.
\]
\note{dans la première ligne, les deux derniers indices du second membre
sont surchargés ; $(a_1, \ldots, a_5)$ du premier membre est souligné}

\note{croquis de la moitié droite de la page : un premier dessin à deux
croix de droites, marqué $a$, $b$, $a'$, $b'$, $c$, entièrement biffé ; à
côté, le même au propre : deux croix superposées, de labels $a$, $b'$,
$a'$, $b$, $c$, $b'$, $a$, $b$, $a'$ ; une étoile de trois droites par un
point, de labels $a'$, $b$, $c$, le côté $a$ en pointillé ; un losange coupé
par sa diagonale $a$ épaissie, faces $f$ et $f'$, côtés $b$, $c$, $b'$,
$c'$ ; en bas, un triangle hachuré dont un côté est épaissi, entouré des
droites $a_1, \ldots, a_5$}

\page{32}

\begin{align*}
  \rho_f = \sigma_1 \sigma_0 &: (a_1, a_2, a_3, a_4, a_5) \longmapsto
    (a_4, a_1, a_3, a_2, a_5) \\
  \rho_s = \sigma_2 \sigma_1 &: (a_1, a_2, a_3, a_4, a_5) \longmapsto
    (a_2, a_3, a_4, a_5, a_1) \\
  \sigma = \sigma_0 \sigma_1 &: (a_1, a_2, a_3, a_4, a_5) \longmapsto
    (a_1, a_3, a_2, a_5, a_4)
\end{align*}
\note{la lettre $\rho$ est lue avec doute ; le reste de la page est vide}

\end{document}
